\documentclass[10pt,a4paper]{article}

% --- PACKAGES ---
\usepackage[utf8]{inputenc}
\usepackage[T1]{fontenc}
\usepackage[scaled=0.92]{helvet}
\usepackage{mathpazo}
\usepackage[margin=1cm, top=1.2cm, bottom=1.2cm]{geometry}
\usepackage{amsmath,amssymb}
\usepackage{booktabs}
\usepackage{tabularx}
\usepackage{multicol}
\usepackage{enumitem}
\usepackage{titlesec}
\usepackage{fancyhdr}
\usepackage{xcolor}
\usepackage{listings}
\usepackage[hidelinks]{hyperref}
\usepackage{setspace}
\usepackage{multicol}

% --- SETUP ---
\setlength{\parindent}{0pt}
\setlength{\parskip}{0.3em}
\setlist[itemize]{leftmargin=1em, noitemsep, topsep=1pt, parsep=0pt, itemsep=0pt}
\setlist[enumerate]{leftmargin=1.2em, noitemsep, topsep=1pt, parsep=0pt, itemsep=0pt}
\setstretch{1.05} % Spasi baris sedikit dilebarkan agar tidak sesak

% --- COLORS & STYLES (MINIMALIST) ---
\definecolor{primary}{HTML}{1F3A5F} % Navy Blue elegan
\definecolor{codebg}{HTML}{F4F5F7}  % Abu-abu sangat muda
\definecolor{codegray}{HTML}{6A737D}

% Section styling
\titleformat{\section}{\normalfont\large\bfseries\sffamily\color{primary}}{}{0em}{}[\vspace{-0.5em}\rule{\linewidth}{0.5pt}]
\titlespacing*{\section}{0pt}{1.2ex}{0.8ex}

% Code Block Styling
\lstset{
    basicstyle=\ttfamily\footnotesize,
    keywordstyle=\color{primary}\bfseries,
    backgroundcolor=\color{codebg},
    frame=none,
    breaklines=true,
    aboveskip=4pt,
    belowskip=4pt,
    xleftmargin=0.5em,
    xrightmargin=0.5em,
    columns=fullflexible,
    keepspaces=true
}

% Header & Footer
\pagestyle{fancy}
\fancyhf{}
\fancyhead[L]{\small\sffamily\bfseries\color{primary} CHEAT SHEET BASIS DATA (P1--P7)}
\fancyhead[R]{\small\sffamily Muhammad Riduwan Khafidi}
\renewcommand{\headrulewidth}{0.4pt}
\renewcommand{\headrule}{\hbox to\headwidth{\color{gray!50}\leaders\hrule height \headrulewidth\hfill}}
\setlength{\headheight}{22pt}

\begin{document}

\vspace{-1em}
\begin{multicols*}{2}

% ==========================================
% 1. KONSEP DASAR & ERD
% ==========================================
\section*{1. Konsep Dasar \& ERD}

\textbf{Kamus Data:} \\
Data = fakta mentah. Database = kumpulan data terorganisir. DBMS = *software* (MySQL, Oracle). RDBMS = DBMS berbasis relasi (tabel).

\textbf{Notasi ERD:} \\
\begin{tabularx}{\linewidth}{@{}l X@{}}
    \textbf{Kotak} & Entitas (Objek nyata) \\
    \textbf{Elips} & Atribut (Properti) \\
    \textbf{Belah Ketupat} & Relasi (Hubungan) \\
    \textbf{Kotak Ganda} & Entitas Lemah (*Weak Entity*) \\
\end{tabularx}

\textbf{Kardinalitas:} 
\begin{itemize}
    \item \textbf{1:1} (Satu Dept punya satu Manajer).
    \item \textbf{1:N} (Satu Dept punya banyak Pegawai).
    \item \textbf{N:M} (Banyak Mhs ambil banyak Matkul).
\end{itemize}

\textbf{EERD (Lanjutan):}
\begin{itemize}
    \item \textbf{Generalization/Specialization:} Sub-grup (Segitiga IS-A). Misal: Entitas \textit{Pegawai} menjadi \textit{Dosen} \& \textit{Staff}.
    \item \textbf{Aggregation:} Sebuah relasi diperlakukan sebagai entitas baru.
\end{itemize}

% ==========================================
% 2. KONVERSI ERD KE TABEL
% ==========================================
\section*{2. Konversi ERD ke Tabel}

Aturan mutlak saat merancang skema relasional:
\begin{enumerate}
    \item \textbf{Entitas Kuat:} Jadi tabel. PK menjadi Primary Key.
    \item \textbf{Entitas Lemah:} Jadi tabel. PK entitas kuat masuk sebagai FK (beri \texttt{ON DELETE CASCADE}).
    \item \textbf{Relasi 1:1:} Pilih salah satu tabel, masukkan PK tabel lain sebagai FK.
    \item \textbf{Relasi 1:N:} PK dari sisi "1" masuk menjadi FK di sisi "N".
    \item \textbf{Relasi N:M:} \textbf{Wajib buat tabel baru} (Tabel Penghubung). Isinya PK dari kedua entitas, jadikan PK bersama (*Composite Key*).
    \item \textbf{Atribut Multivalued:} Buat tabel baru (PK entitas asli + atribut tersebut).
\end{enumerate}

% ==========================================
% 3. DDL (STRUKTUR)
% ==========================================
\section*{3. DDL (Data Definition Language)}

\textbf{Tipe Data Kunci:}
\begin{itemize}
    \item \texttt{CHAR(n)}: Panjang tetap (boros, cepat).
    \item \texttt{VARCHAR(n)}: Panjang dinamis (hemat).
    \item \texttt{INT}, \texttt{FLOAT}, \texttt{NUMERIC(p,d)}, \texttt{DATE}, \texttt{NULL}.
\end{itemize}

\textbf{Create Table \& Constraints:}
\begin{lstlisting}[language=SQL]
CREATE TABLE branch (
  branch_name VARCHAR(15) NOT NULL,
  assets INT CHECK (assets >= 0),
  PRIMARY KEY (branch_name)
);

CREATE TABLE borrower (
  cust_name VARCHAR(30),
  loan_num INT,
  PRIMARY KEY (cust_name, loan_num),
  FOREIGN KEY (cust_name) 
    REFERENCES customer(cust_name)
    ON DELETE CASCADE
);
\end{lstlisting}

\textbf{Alter Table:}
\begin{lstlisting}[language=SQL]
ALTER TABLE borrower ADD b_date DATE;
ALTER TABLE borrower CHANGE b_date birth_date DATE;
ALTER TABLE borrower DROP birth_date;
\end{lstlisting}

\textbf{Index (Pencarian Cepat):}
\begin{lstlisting}[language=SQL]
CREATE INDEX idx_name ON Account (branch_name);
\end{lstlisting}
\textit{Catatan: Jangan buat index di tabel kecil atau tabel yang sering di-Update.}

% ==========================================
% 4. DML (MANIPULASI DATA)
% ==========================================
\section*{4. DML (Data Manipulation)}

\textbf{Struktur Dasar:}
\begin{lstlisting}[language=SQL]
SELECT kolom FROM tabel WHERE kondisi;
\end{lstlisting}

\textbf{Operator WHERE (Filter):}
\begin{itemize}
    \item Logika: \texttt{AND, OR, NOT}.
    \item Perbandingan: \texttt{=, >, >=, <, <=, <>}.
    \item Teks: \texttt{LIKE} (\% = banyak huruf, \_ = 1 huruf). \\
    Contoh: \texttt{LIKE 'Fr\%'} (awalan Fr), \texttt{LIKE '\_\_re'} (2 huruf + re).
    \item Kosong: \texttt{IS NULL} / \texttt{IS NOT NULL}.
    \item Rentang: \texttt{BETWEEN 10 AND 20} / \texttt{IN ('102', '104')}.
\end{itemize}

\textbf{Fungsi Agregat:} \texttt{AVG, MIN, MAX, SUM, COUNT}.
\begin{lstlisting}[language=SQL]
SELECT COUNT(*), AVG(Balance) FROM Account;
\end{lstlisting}

\textbf{Join 2 Tabel (Manual):} \\
Gunakan \texttt{WHERE} untuk menghubungkan PK dan FK agar tidak terjadi \textit{Cartesian Product}.
\begin{lstlisting}[language=SQL]
SELECT A.Owner, D.Amount
FROM Account A, Deposit D
WHERE A.Number = D.Account 
  AND D.Amount > 1000;
\end{lstlisting}

\textbf{CRUD Dasar:}
\begin{lstlisting}[language=SQL]
INSERT INTO Account VALUES ('A1', 5000);
UPDATE Account SET balance = balance * 1.05;
DELETE FROM Account WHERE balance < 100;
\end{lstlisting}

\textbf{CASE Statement (If-Else SQL):}
\begin{lstlisting}[language=SQL]
UPDATE Account SET balance = 
  CASE 
    WHEN balance > 1000 THEN balance * 1.06
    ELSE balance * 1.05 
  END;
\end{lstlisting}

% ==========================================
% 5. TIPS & TRIK (CONTEKAN)
% ==========================================
\section*{Trik Ujian \& Praktikum}

\noindent\textcolor{primary}{\textbf{1. Jebakan NULL}} \\
Jangan pernah tulis \texttt{= NULL}. SQL tidak mengenali itu. \textbf{Selalu} gunakan \texttt{IS NULL}.

\vspace{0.5em}
\noindent\textcolor{primary}{\textbf{2. Aturan N:M}} \\
Relasi Many-to-Many (N:M) \textbf{TIDAK BISA} dimasukkan sebagai FK di tabel yang sudah ada. Wajib bikin tabel baru.

\vspace{0.5em}
\noindent\textcolor{primary}{\textbf{3. Bahaya Cartesian}} \\
\texttt{SELECT * FROM A, B} tanpa \texttt{WHERE} = Bencana (Data berlipat ganda tak beraturan). Selalu hubungkan PK-FK.

\vspace{0.5em}
\noindent\textcolor{primary}{\textbf{4. DELETE vs DROP}} \\
\texttt{DELETE} = Menghapus baris data (Tabel tetap ada). \\
\texttt{DROP} = Menghancurkan struktur tabel beserta datanya (Tabel hilang).

\vspace{0.5em}
\noindent\textcolor{primary}{\textbf{5. Cek CHAR vs VARCHAR}} \\
Gunakan \texttt{CHAR} jika data dipastikan panjangnya sama (misal Kode Negara 'ID'). Gunakan \texttt{VARCHAR} untuk Nama/Alamat.

\vspace{0.5em}
\noindent\textcolor{primary}{\textbf{6. FK \& CASCADE}} \\
Jika induk dihapus, tapi anak masih pakai datanya, SQL akan *error*. Solusinya tambahkan \texttt{ON DELETE CASCADE} saat membuat FK.

\end{multicols*}
\end{document}